\textbf{Phần 1. Trắc nghiệm lựa chọn}
\setcounterpageref{ex}{0}

\begin{ex} %[4D8N40]
Tính chất giao hoán của phép nhân là
\choice
{$a \times b = a + b$}
{$(a \times b) \times c = a \times (b \times c)$}
{\True $a \times b = b \times a$}
{$a \times 1 = a$}
\loigiai{
$a \times b = b \times a$.\\
\textbf{Đáp án của câu hỏi là C}
}
\end{ex}

\begin{ex} %[4D8N40]
Tính chất kết hợp của phép nhân là
\choice
{$a \times b = b \times a$}
{\True $(a \times b) \times c = a \times (b \times c)$}
{$a \times (b + c) = a \times b + a \times c$}
{$a \times 0 = 0$}
\loigiai{
$(a \times b) \times c = a \times (b \times c)$.\\
\textbf{Đáp án của câu hỏi là B}
}
\end{ex}

\begin{ex} %[4D8H40]
Tính thuận tiện: $25 \times 37 \times 4$.
\choice
{$3600$}
{$3800$}
{\True $3700$}
{$3500$}
\loigiai{
$(25 \times 4) \times 37 = 100 \times 37 = 3700$.\\
\textbf{Đáp án của câu hỏi là C}
}
\end{ex}

\begin{ex} %[4D8H40]
Tính thuận tiện: $125 \times 48 \times 8$.
\choice
{$46\,000$}
{\True $48\,000$}
{$50\,000$}
{$480\,000$}
\loigiai{
$(125 \times 8) \times 48 = 1000 \times 48 = 48\,000$.\\
\textbf{Đáp án của câu hỏi là B}
}
\end{ex}

\begin{ex} %[4D8H40]
Tính thuận tiện: $50 \times 19 \times 2$.
\choice
{\True $1900$}
{$1800$}
{$2000$}
{$1950$}
\loigiai{
$(50 \times 2) \times 19 = 100 \times 19 = 1900$.\\
\textbf{Đáp án của câu hỏi là A}
}
\end{ex}

\begin{ex} %[4D8H40]
Tính thuận tiện: $2 \times 456 \times 5$.
\choice
{$4460$}
{\True $4560$}
{$4660$}
{$4500$}
\loigiai{
$(2 \times 5) \times 456 = 10 \times 456 = 4560$.\\
\textbf{Đáp án của câu hỏi là B}
}
\end{ex}

\begin{ex} %[4D8V40]
Tính nhanh giá trị biểu thức: $4 \times 18 \times 25$.
\choice
{$1600$}
{$2000$}
{$1700$}
{\True $1800$}
\loigiai{
$(4 \times 25) \times 18 = 100 \times 18 = 1800$.\\
\textbf{Đáp án của câu hỏi là D}
}
\end{ex}

\begin{ex} %[4D8V40]
Tính nhanh: $8 \times 35 \times 125$.
\choice
{\True $35\,000$}
{$34\,000$}
{$36\,000$}
{$350\,000$}
\loigiai{
$(8 \times 125) \times 35 = 1000 \times 35 = 35\,000$.\\
\textbf{Đáp án của câu hỏi là A}
}
\end{ex}

\begin{ex} %[4D8V40]
Một phòng học có 8 dãy bàn, mỗi dãy có 5 bàn, mỗi bàn có 2 học sinh ngồi. Số học sinh trong phòng học đó tính thuận tiện là
\choice
{$70$ học sinh}
{$90$ học sinh}
{$100$ học sinh}
{\True $80$ học sinh}
\loigiai{
$8 \times (5 \times 2) = 8 \times 10 = 80$ học sinh.\\
\textbf{Đáp án của câu hỏi là D}
}
\end{ex}

\begin{ex} %[4D8V40]
Một cửa hàng xếp 4 thùng bánh, mỗi thùng có 25 hộp bánh, mỗi hộp bánh giá $15\,000$ đồng. Tổng giá trị của 4 thùng bánh là
\choice
{\True $1\,500\,000$ đồng}
{$1\,400\,000$ đồng}
{$1\,600\,000$ đồng}
{$1\,800\,000$ đồng}
\loigiai{
$(4 \times 25) \times 15\,000 = 100 \times 15\,000 = 1\,500\,000$ đồng.\\
\textbf{Đáp án của câu hỏi là A}
}
\end{ex}

\begin{ex} %[4D8V40]
Tìm giá trị của $m$ biết: $m \times 25 \times 4 = 8400$.
\choice
{$82$}
{$86$}
{$88$}
{\True $84$}
\loigiai{
$m \times 100 = 8400 \Rightarrow m = 84$.\\
\textbf{Đáp án của câu hỏi là D}
}
\end{ex}

\begin{ex} %[4D8C40]
Tính tích: $A = 2 \times 4 \times 5 \times 25 \times 8 \times 125$.
\choice
{$1\,000\,000$}
{$10\,000\,000$}
{\True $1\,000\,000$}
{$100\,000$}
\loigiai{
$(2 \times 5) \times (4 \times 25) \times (8 \times 125) = 10 \times 100 \times 1000 = 1\,000\,000$.\\
\textbf{Đáp án của câu hỏi là C}
}
\end{ex}

\textbf{Phần 2. Trắc nghiệm đúng sai}
\setcounterpageref{ex}{0}

\begin{ex} %[4D8H40]
Xét tính đúng sai về tính chất phép nhân:
\choiceTF
{\True $a \times b = b \times a$.}
{$(a \times b) \times c = a \times b + a \times c$.}
{\True $(a \times b) \times c = a \times (b \times c)$.}
{\True $25 \times 18 \times 4 = (25 \times 4) \times 18 = 1800$.}
\loigiai{
\begin{itemize}
    \item Giao hoán: Đúng.
    \item Đây không phải tính chất kết hợp. Sai.
    \item Kết hợp: Đúng.
    \item Tính thuận tiện đúng.
\end{itemize}
\begin{itemize}
    \item a) \textbf{Đúng}.
    \item b) \textbf{Sai}.
    \item c) \textbf{Đúng}.
    \item d) \textbf{Đúng}.
\end{itemize}
}
\end{ex}

\begin{ex} %[4D8V40]
Cho hai biểu thức: $M = 125 \times 36 \times 8$ và $N = 25 \times 36 \times 40$. Xét tính đúng sai:
\choiceTF
{\True $M = 36\,000$.}
{\True $N = 36\,000$.}
{Giá trị biểu thức $M$ lớn hơn giá trị biểu thức $N$.}
{\True $M = N$.}
\loigiai{
\begin{itemize}
    \item $M = (125 \times 8) \times 36 = 1000 \times 36 = 36\,000$. Đúng.
    \item $N = (25 \times 40) \times 36 = 1000 \times 36 = 36\,000$. Đúng.
    \item $M = N$ nên $M > N$ là Sai.
    \item $M = N$. Đúng.
\end{itemize}
\begin{itemize}
    \item a) \textbf{Đúng}.
    \item b) \textbf{Đúng}.
    \item c) \textbf{Sai}.
    \item d) \textbf{Đúng}.
\end{itemize}
}
\end{ex}

\textbf{Phần 3. Trả lời ngắn}
\setcounterpageref{ex}{0}

\begin{ex} %[4D8H40]
Tính thuận tiện: $25 \times 49 \times 4$.
\shortans{4900}
\loigiai{
\textbf{Đáp số:} $4900$\\
$(25 \times 4) \times 49 = 4900$.
}
\end{ex}

\begin{ex} %[4D8H40]
Tính thuận tiện: $125 \times 73 \times 8$.
\shortans{73000}
\loigiai{
\textbf{Đáp số:} $73000$\\
$(125 \times 8) \times 73 = 73\,000$.
}
\end{ex}

\begin{ex} %[4D8V40]
Tính nhanh: $50 \times 68 \times 2$.
\shortans{6800}
\loigiai{
\textbf{Đáp số:} $6800$\\
$(50 \times 2) \times 68 = 6800$.
}
\end{ex}

\begin{ex} %[4D8V40]
Tìm $x$ biết: $x : (25 \times 4) = 45$.
\shortans{4500}
\loigiai{
\textbf{Đáp số:} $4500$\\
$x : 100 = 45 \Rightarrow x = 4500$.
}
\end{ex}

\begin{ex} %[4D8V40]
Một đoàn xe có 5 xe, mỗi xe chở 20 bao gạo, mỗi bao gạo nặng $50\text{ kg}$. Tính tổng khối lượng gạo đoàn xe chở theo tấn.
\shortans{5}
\loigiai{
\textbf{Đáp số:} $5$\\
$5 \times 20 \times 50 = 5000\text{ kg} = 5$ tấn.
}
\end{ex}

\begin{ex} %[4D8C40]
Tính tích: $2 \times 5 \times 4 \times 25 \times 8 \times 125$.
\shortans{1000000}
\loigiai{
\textbf{Đáp số:} $1000000$\\
$10 \times 100 \times 1000 = 1\,000\,000$.
}
\end{ex}
